\begin{abstract}
Reference outcomes are costly, while cheaper surrogates are available for most
units. We study the conditional distribution of a reference outcome at a
chosen covariate value under selective verification, where few verified
observations may lie nearby. A surrogate-free \(X\)-based estimator remains
valid but can be noisy; surrogate outcome regression improves precision but can
increase error when its local correction is inaccurate. We retain this
surrogate-free estimator as an anchor and select one correction weight across
the CDF by an observable
local variance criterion. Sample splitting separates regression fitting and
weight selection from final CDF evaluation. We establish path-oracle
equivalence, a uniform Gaussian limit, and multiplier validity when the
verification fraction may decrease. The bands cover a kernel-localized CDF;
point-target CDF and quantile inference additionally require localization-bias
control. Simulations show that the criterion distinguishes helpful from harmful
full borrowing: the adaptive weight approaches zero under local reversal and
retains gains under alignment, although Full can be more efficient under strong
alignment. Two real-data applications and a controlled review benchmark show
that the precision gain varies with the strength and direction of the local
surrogate relation.
\end{abstract}

\noindent\textbf{Keywords:} two-phase sampling; missing at random; negative transfer; selective verification; semiparametric efficiency; multiplier bootstrap.
